% Generated draft: requires mathematical and source review.
\documentclass[11pt]{article}
\usepackage[margin=25mm]{geometry}
\usepackage{fontspec}
\setmainfont{Latin Modern Roman}
\usepackage{amsmath,amssymb,mathtools,unicode-math}
\setmathfont{Latin Modern Math}
\usepackage{enumitem,needspace,xurl,hyperref}
\setlength{\parindent}{0pt}
\setlength{\parskip}{5pt}
\setlength{\emergencystretch}{3em}
\newcommand{\N}{\mathbb{N}}
\newcommand{\Z}{\mathbb{Z}}
\newcommand{\Q}{\mathbb{Q}}
\newcommand{\R}{\mathbb{R}}
\newcommand{\C}{\mathbb{C}}
\newcommand{\F}{\mathbb{F}}
\newcommand{\A}{\mathbb{A}}
\newcommand{\PP}{\mathbb{P}}
\newcommand{\E}{\mathbb{E}}
\newcommand{\eps}{\varepsilon}
\newcommand{\dd}{\,\mathrm{d}}
\begin{document}
\section*{UnsolvedMath ID: 10900042}
\textbf{Source problem code:} catalogue entry
\subsection{Definitions and mathematical statement}
Identify \(S^2\) with the ideal boundary of \(\mathbb H^3\). A Kleinian group is a discrete subgroup \(\Gamma\subset\mathrm{PSL}_2(\C)\), acting by orientation-preserving isometries of \(\mathbb H^3\). Its limit set is
\[
 \Lambda(\Gamma)=\overline{\Gamma\cdot o}\cap S^2,
\]
where closure is in \(\mathbb H^3\cup S^2\) and \(o\in\mathbb H^3\); the result is independent of \(o\). In the usual non-elementary setting, \(\Gamma\) is convex cocompact when the quotient of the hyperbolic convex hull of \(\Lambda(\Gamma)\) by \(\Gamma\) is compact.

Walsh's Problem 5.17 asks to characterize the subsets \(A\subset S^2\) for which there exists a convex-cocompact Kleinian group with
\[
 A=\Lambda(\Gamma).
\]
The source does not specify whether subsets are classified literally in the given sphere, up to ambient homeomorphism, up to a Möbius transformation, or up to quasiconformal equivalence. These are different classification targets. Torsion-free versus orbifold groups and elementary versus non-elementary groups are also not specified.

The supplied summary names Cantor sets, Jordan curves and Sierpiński carpets and links the full question to unresolved carpet questions. It supplies no hypotheses making that list exhaustive or establishing the claimed equivalence. For example, the standard definition also allows cocompact lattice cases with full-sphere limit set unless they are explicitly excluded. Thus the summary's list and equivalence cannot replace the full subset-classification request; they remain qualified context requiring source review.
\subsection{Short English statement}
Characterize the subsets of the sphere at infinity that occur as limit sets of convex-cocompact Kleinian groups. The group class and equivalence relation on subsets need specification. The supplied named topological examples and carpet discussion do not constitute an exhaustive verified classification.
\subsection{Sources}
\begin{itemize}
\item Local UnsolvedMath registry, permanent ID \(10900042\), exact source code \texttt{catalogue entry}, subsection 5.17, attributed to Walsh; no proposed year supplied.
\item Delp, Hoffoss and Manning, \emph{Problems In Groups, Geometry, and Three-Manifolds} (2015), 5.17, PDF page 5: \url{https://arxiv.org/pdf/1512.04620#page=5}; supplied status-source URL \url{https://arxiv.org/abs/1512.04620}. No source-side resolution is recorded.
\item The research summary names Cantor/Jordan/Sierpiński-carpet phenomena and a claimed relation to unresolved carpet questions, without complete references or the needed hypotheses. It says no single verified theorem settles the full classification.
\item Source-TeX extraction checked against rendered PDF text, accessed 2026-07-29. Source evidence NEEDS\_REVIEW: open in 2015, no complete modern primary-literature audit recorded. Registry status open and research classification OPEN-TRIAGE.
\item Dataset attribution: Ulam AI / UnsolvedMath Contributors, CC BY 4.0. Source rights NEEDS\_REVIEW: publicly accessible author-submitted source TeX and PDF, redistribution terms requiring release review.
\end{itemize}
\end{document}
